\begin{resumen}
	Esta tesis aborda el problema de recuperar las expresiones analíticas de todas las soluciones de un sistema de ecuaciones no lineales paramétrico. A partir de una cuadrícula de parámetros generada mediante producto cartesiano, se resuelve numéricamente el sistema para cada tupla utilizando múltiples puntos iniciales, construyendo así un conjunto de datos que contiene todas las soluciones posibles. La principal dificultad reside en que para un mismo valor de los parámetros pueden coexistir varias soluciones distintas, cuyos puntos se entremezclan en el conjunto de entrenamiento. Se propone un algoritmo iterativo de regresión simbólica que descubre secuencialmente cada una de las funciones solución, retirando los puntos ya explicados antes de la siguiente iteración. Se comparan tres funciones de pérdida (error cuadrático medio, conteo de coincidencias y pérdida sigmoidal) sobre 30 ecuaciones univariadas, seleccionando la sigmoidal por su mayor tasa de éxito (90\,\%) y menor tiempo de cómputo. La metodología completa se evalúa sobre 30 sistemas de dos variables y hasta tres parámetros, logrando recuperar todas las soluciones en el 90\,\% de los casos, incluyendo sistemas con tres y cuatro expresiones diferentes para una misma tupla de parámetros. Los resultados demuestran que la combinación de la búsqueda iterativa con la pérdida sigmoidal permite obtener expresiones analíticas válidas y de alta precisión, facilitando el análisis de sensibilidad y la predicción del comportamiento del sistema sin necesidad de resolverlo numéricamente cada vez.

	


	\textbf{Palabras clave:} Sistemas de ecuaciones no lineales paramétricos, regresión simbólica, programación genética, descubrimiento de funciones, solución múltiple, pérdida sigmoidal, PySR.
\end{resumen}

\begin{abstract}
	This thesis addresses the problem of retrieving analytical expressions for all solutions of a parametric system of nonlinear equations. Starting from a grid of parameters generated through a Cartesian product, the system is solved numerically for each tuple using multiple initial points, thereby constructing a dataset containing all possible solutions. The main challenge lies in the fact that for the same parameter values, multiple distinct solutions may coexist, with their points intermingling in the training set. An iterative symbolic regression algorithm is proposed that sequentially discovers each solution function, removing already explained points before the next iteration. Three loss functions (mean squared error, match count, and sigmoidal loss) are compared across 30 univariate equations, with the sigmoidal function selected for its higher success rate (90\,\%) and lower computational time. The complete methodology is evaluated on 30 two-variable systems with up to three parameters, successfully recovering all solutions in 90\,\% of cases, including systems with three and four different expressions for a given parameter tuple. The results demonstrate that combining iterative search with sigmoidal loss enables obtaining valid and highly accurate analytical expressions, facilitating sensitivity analysis and behavior prediction of the system without needing to solve it numerically each time.

	

	
	\textbf{Keywords:} Parametric nonlinear equation systems, symbolic regression, genetic programming, function discovery, multiple solutions, sigmoidal loss, PySR.
\end{abstract}